\chapter{Pränexierung und QCIR-Repräsentation der
\emph{preferred}-Kodierung}
\label{ch:preferred-prenex-qcir}

Kapitel~\ref{ch:semantic-independence} charakterisiert die bedingte
Unabhängigkeit durch die QBF
\(\Phi_F^\sigma(X,Y\mid Z)\). Für die algorithmische Auswertung muss
diese mathematische Darstellung in eine konkrete Solverrepräsentation
überführt werden. Der vorliegende Anhang beschreibt diesen Schritt für
die \emph{preferred}-Semantik.

Dabei sind zwei Transformationen zu unterscheiden. Zunächst werden die
innerhalb der \emph{preferred}-Prädikate auftretenden Quantoren in einen
gemeinsamen Quantorenpräfix überführt. Anschließend wird die
quantorenfreie Matrix als boolescher Schaltkreis im Format
\texttt{QCIR-G14} dargestellt.

Die resultierende Formel besitzt somit Pränexform, ihre Matrix liegt
jedoch nicht in konjunktiver Normalform vor. Es handelt sich insbesondere
nicht um eine PCNF- beziehungsweise QDIMACS-Kodierung. Eine solche
Konvertierung ist für die verwendete Werkzeugkette nicht erforderlich:
Wie in Kapitel~\ref{ch:empirical-analysis} beschrieben, werden die
erzeugten \texttt{QCIR-G14}-Dateien ohne weitere Formatumwandlung an
QuAbS übergeben. Obwohl \texttt{QCIR-G14} auch nicht-pränexe Formeln
darstellen kann, erzeugt der verwendete Übersetzer einen vollständig
pränexen Quantorenpräfix mit einer nicht-KNF-Schaltkreismatrix.

\section{Überführung der \emph{preferred}-QBF in Pränexform}
\label{sec:preferred-prenex}

Für \(\sigma=\mathsf{pr}\) ergibt sich aus
Definition~\ref{def:generic-independence-qbf} zunächst
\[
\begin{split}
\Phi_F^{\mathsf{pr}}(X,Y\mid Z)
:=
\forall\mathbf V^1\,\forall\mathbf V^2\,
\Bigl(&
    \operatorname{Pr}_F(\mathbf V^1)
    \land
    \operatorname{Pr}_F(\mathbf V^2)
    \land
    \operatorname{Eq}^{1,2}_Z
\\
&\rightarrow
    \exists\mathbf V^3\,
    \bigl(
        \operatorname{Pr}_F(\mathbf V^3)
        \land
        \operatorname{Mix}^{1,2,3}_{X,Y\mid Z}
    \bigr)
\Bigr).
\end{split}
\]

Vor der Pränexierung werden die gebundenen Variablen der drei
\emph{preferred}-Prädikate durch paarweise disjunkte Blöcke
\(\widehat{\mathbf V}^{1}\),
\(\widehat{\mathbf V}^{2}\) und
\(\widehat{\mathbf V}^{3}\) ersetzt. Für \(k\in\{1,2,3\}\) sei
\[
\operatorname{Max}_F^k
:=
\neg\Bigl(
    \operatorname{Co}_F(\widehat{\mathbf V}^{k})
    \land
    I^k<\widehat{I}^{k}
\Bigr).
\]
Damit gilt
\[
\operatorname{Pr}_F(\mathbf V^k)
=
\operatorname{Co}_F(\mathbf V^k)
\land
\forall\widehat{\mathbf V}^{k}\,
\operatorname{Max}_F^k.
\]

Für die spätere Matrix werden die beiden quantorenfreien Teilformeln
\[
\begin{aligned}
\mathcal A_F
:={}&
    \operatorname{Co}_F(\mathbf V^1)
    \land
    \operatorname{Co}_F(\mathbf V^2)
    \land
    \operatorname{Eq}^{1,2}_Z
    \land
    \operatorname{Max}_F^1
    \land
    \operatorname{Max}_F^2,
\\
\mathcal C_F
:={}&
    \operatorname{Co}_F(\mathbf V^3)
    \land
    \operatorname{Mix}^{1,2,3}_{X,Y\mid Z}
    \land
    \operatorname{Max}_F^3
\end{aligned}
\]
verwendet.

\begin{Proposition}
\label{prop:preferred-prenex-equivalent}
Die QBF \(\Phi_F^{\mathsf{pr}}(X,Y\mid Z)\) ist äquivalent zu der
pränexen QBF
\[
\begin{split}
\forall\mathbf V^1\,
\forall\mathbf V^2\,
\exists\mathbf V^3\,
\exists\widehat{\mathbf V}^{1}\,
\exists\widehat{\mathbf V}^{2}\,
\forall\widehat{\mathbf V}^{3}\,
\bigl(
    \mathcal A_F\rightarrow\mathcal C_F
\bigr).
\end{split}
\]
Zusammengefasst besitzt ihr Präfix die Alternationsstruktur
\[
\forall(\mathbf V^1\cup\mathbf V^2)\,
\exists(\mathbf V^3\cup
        \widehat{\mathbf V}^{1}\cup
        \widehat{\mathbf V}^{2})\,
\forall\widehat{\mathbf V}^{3}.
\]
\end{Proposition}

\begin{proof}
Nach der Umbenennung der gebundenen Variablen sind alle sechs
Labellingvariablenblöcke paarweise disjunkt. Für eine Variable
beziehungsweise einen Variablenblock \(\mathbf W\), der nicht frei in
\(\beta\) vorkommt, gilt
\[
(\forall\mathbf W\,\alpha)\rightarrow\beta
\equiv
\exists\mathbf W\,(\alpha\rightarrow\beta).
\]
Daher wechseln die Quantoren über
\(\widehat{\mathbf V}^{1}\) und
\(\widehat{\mathbf V}^{2}\) beim Herausziehen aus dem Antezedens ihre
Polarität: Aus den universellen Maximalitätsprüfungen innerhalb der
beiden \emph{preferred}-Prädikate werden existentielle Blöcke im
gemeinsamen Präfix.

Für den existentiellen Block des rekombinierenden Labellings gilt
dagegen
\[
\alpha\rightarrow\exists\mathbf W\,\beta
\equiv
\exists\mathbf W\,(\alpha\rightarrow\beta),
\]
sofern \(\mathbf W\) nicht frei in \(\alpha\) vorkommt. Der Block
\(\mathbf V^3\) bleibt daher existentiell. Die Maximalitätsprüfung von
\(\mathbf V^3\) steht im Konsequens und behält ihren universellen
Quantor über \(\widehat{\mathbf V}^{3}\).

Die Blöcke \(\mathbf V^3\),
\(\widehat{\mathbf V}^{1}\) und
\(\widehat{\mathbf V}^{2}\) sind sämtlich existentiell und können daher
innerhalb des mittleren Quantorenblocks vertauscht werden. Die
angegebene Reihenfolge entspricht der verwendeten Implementierung.
\end{proof}

Nach Eliminierung der letzten Implikation ist die quantorenfreie Matrix
\[
\Psi_F
:=
\neg\mathcal A_F\lor\mathcal C_F.
\]
Diese Matrix muss für die Übergabe an QuAbS nicht in KNF umgewandelt
werden. Sie wird im nächsten Abschnitt unmittelbar als boolescher
Schaltkreis dargestellt.

Jeder der sechs Labellingvariablenblöcke enthält für jedes Argument die
drei Variablen für \(\mathsf{in}\), \(\mathsf{out}\) und
\(\mathsf{und}\). Für \(n=|A|\) besitzt die pränexe QBF daher insgesamt
\[
6\cdot 3n=18n
\]
quantifizierte Variablen.

\begin{Example}[Zweizyklus]
\label{ex:preferred-prenex-two-cycle}
Sei
\[
F_{\leftrightarrow}
=
\bigl(
    \{a_1,a_2\},
    \{(a_1,a_2),(a_2,a_1)\}
\bigr)
\]
und seien
\[
X=\{a_1\},
\qquad
Y=\{a_2\},
\qquad
Z=\emptyset.
\]
Der AF besitzt die beiden \emph{preferred}-Labellings
\[
\begin{array}{c|ccc}
    & \mathsf{in} & \mathsf{out} & \mathsf{und}\\
    \midrule
    \lambda_1 & \{a_1\} & \{a_2\} & \emptyset\\
    \lambda_2 & \{a_2\} & \{a_1\} & \emptyset
\end{array}
\]
sowie ein nicht maximales \emph{complete}-Labelling, in dem beide
Argumente mit \(\mathsf{und}\) bezeichnet sind.

Für dieses nicht maximale Labelling kann ein existentieller
Vergleichsblock, beispielsweise
\(\widehat{\mathbf V}^{1}\), das Labelling \(\lambda_1\) kodieren. Dann
gilt
\[
\emptyset
<
\mathsf{in}(\lambda_1)=\{a_1\},
\]
wodurch \(\operatorname{Max}_F^1\) falsch und damit das Antezedens der
Implikation deaktiviert wird. Auf diese Weise werden Belegungen, die
keine \emph{preferred}-Labellings kodieren, korrekt von der
Rekombinationsanforderung ausgeschlossen.

Kodieren \(\mathbf V^1\) und \(\mathbf V^2\) dagegen
\(\lambda_1\) beziehungsweise \(\lambda_2\), so verlangt
\(\operatorname{Mix}^{1,2,3}_{X,Y\mid Z}\), dass das dritte Labelling
sowohl \(a_1\) als auch \(a_2\) mit \(\mathsf{in}\) bezeichnet. Ein
solches \emph{complete}- und damit auch kein solches
\emph{preferred}-Labelling existiert. Folglich ist
\[
\Phi_{F_{\leftrightarrow}}^{\mathsf{pr}}
    (\{a_1\},\{a_2\}\mid\emptyset)
\]
falsch. Für \(n=2\) enthält die pränexe Formel \(36\) quantifizierte
Variablen und besitzt die Blockstruktur
\[
\underbrace{\forall\,12}_{\mathbf V^1,\mathbf V^2}\;
\underbrace{\exists\,18}_{\mathbf V^3,
    \widehat{\mathbf V}^{1},\widehat{\mathbf V}^{2}}\;
\underbrace{\forall\,6}_{\widehat{\mathbf V}^{3}}.
\]
\end{Example}

\section{Übersetzung in das QCIR-G14-Format}
\label{sec:preferred-qcir}

QuAbS verarbeitet die QBF als booleschen Schaltkreis im Format
\texttt{QCIR-G14}. Quantifizierte Variablen und Gatter werden dabei
durch positive ganzzahlige Bezeichner repräsentiert; ein negatives
Vorzeichen bezeichnet die Negation eines Variablen- oder
Gatterausgangs. Konjunktionen und Disjunktionen werden durch
\texttt{and}- beziehungsweise \texttt{or}-Gatter dargestellt.

\subsection{Variablen- und Gatterbelegung}

Die Implementierung verwendet sechs Labellingkopien:
\[
\mathbf V^1,\quad
\mathbf V^2,\quad
\mathbf V^3,\quad
\widehat{\mathbf V}^{1},\quad
\widehat{\mathbf V}^{2},\quad
\widehat{\mathbf V}^{3}.
\]
Sie werden intern als Kopien \(k=1,\ldots,6\) geführt. Für
\(A=\{1,\ldots,n\}\) lautet die Variablenzuordnung
\[
\begin{aligned}
\operatorname{id}(i_a^k)
    &=(k-1)3n+a,\\
\operatorname{id}(o_a^k)
    &=(k-1)3n+n+a,\\
\operatorname{id}(u_a^k)
    &=(k-1)3n+2n+a.
\end{aligned}
\]
Innerhalb jeder Kopie werden somit zuerst alle
\(\mathsf{in}\)-Variablen, anschließend alle
\(\mathsf{out}\)-Variablen und zuletzt alle
\(\mathsf{und}\)-Variablen ausgegeben. Die IDs \(1,\ldots,18n\) sind
quantifizierten Variablen vorbehalten; das erste Gatter erhält die ID
\(18n+1\).

Der Quantorenpräfix wird blockweise durch
\texttt{forall(...)} und \texttt{exists(...)} dargestellt. Aufeinander
folgende gleichartige Zeilen gehören zum selben Quantorenblock. Danach
bestimmt
\[
\texttt{output(}g_{\Phi}\texttt{)}
\]
das Ausgabegatter der gesamten QBF.

Jede zusammengesetzte Teilformel erhält ein eigenes Gatter. Für
Literale \(\ell_1,\ldots,\ell_m\) werden insbesondere
\[
\begin{array}{rcl}
\ell_1\land\cdots\land\ell_m
    &\longmapsto&
    \texttt{g = and(}\ell_1\texttt{, \ldots, }\ell_m\texttt{)},\\
\ell_1\lor\cdots\lor\ell_m
    &\longmapsto&
    \texttt{g = or(}\ell_1\texttt{, \ldots, }\ell_m\texttt{)}
\end{array}
\]
verwendet. Eine Äquivalenz wird ausschließlich durch
\texttt{and}- und \texttt{or}-Gatter ausgedrückt:
\[
x\leftrightarrow y
\equiv
(\neg x\lor y)\land(x\lor\neg y).
\]
Die Maximalitätsbedingung
\[
\neg\bigl(
    \operatorname{Co}_F(\widehat{\mathbf V}^{k})
    \land I^k<\widehat I^k
\bigr)
\]
wird daher durch
\[
\texttt{g\_max = or(-g\_cohat, -g\_lt)}
\]
repräsentiert. Für die äußere Implikation
\(\mathcal A_F\rightarrow\mathcal C_F\) entsteht schließlich das
Ausgabegatter
\[
\texttt{g\_phi = or(-g\_A, g\_C)}.
\]
Die leere Konjunktion \texttt{and()} repräsentiert \(\top\). Dies tritt
beispielsweise bei
\(\operatorname{Eq}^{j,k}_{\emptyset}\) sowie bei der
\emph{complete}-Kodierung eines Arguments ohne Angreifer auf.

\paragraph{Minimales Syntaxbeispiel.}
Die wahre QBF
\[
\forall x\,\exists y\,
\bigl((x\lor\neg y)\land(\neg x\lor y)\bigr)
\]
erhält die folgende vollständige QCIR-Darstellung:
\begin{lstlisting}[
    language={},
    caption={Minimale pränexe QCIR-G14-Eingabe},
    label={lst:qcir-minimal}
]
#QCIR-G14 5
forall(1)
exists(2)
output(5)
3 = or(1, -2)
4 = or(-1, 2)
5 = and(3, 4)
\end{lstlisting}
Für jede Belegung von \(x\) kann \(y=x\) gewählt werden. QuAbS gibt
daher \texttt{r SAT} aus. Die Zahl hinter \texttt{\#QCIR-G14} gibt die
höchste verwendete ID an. QuAbS akzeptiert außerdem die vom
vorliegenden Übersetzer erzeugte Kurzform \texttt{\#QCIR-G14}.

\subsection{Übersetzung des Zweizyklen-Beispiels}

Der AF aus Beispiel~\ref{ex:preferred-prenex-two-cycle} wird im
verwendeten ICCMA-Format durch
\begin{lstlisting}[
    language={},
    caption={ICCMA-Darstellung des Zweizyklus},
    label={lst:iccma-preferred-two-cycle}
]
p af 2
# a1
# a2
1 2
2 1
\end{lstlisting}
beschrieben. Die beiden letzten Zeilen repräsentieren die Angriffe
\(a_1\rightarrow a_2\) und \(a_2\rightarrow a_1\).

Für \(X=\{a_1\}\), \(Y=\{a_2\}\) und \(Z=\emptyset\) kann die
QCIR-Datei durch
\begin{lstlisting}[language={}]
python3 af_to_prind_qcir.py zweizyklus.af -x 1 -y 2 \
    -o zweizyklus.qcir
\end{lstlisting}
erzeugt werden. Da \(n=2\) gilt, belegt jede Labellingkopie sechs
Variablen. Der erzeugte Quantorenpräfix lautet:
\begin{lstlisting}[
    language={},
    caption={QCIR-Präfix des Zweizyklen-Beispiels},
    label={lst:qcir-prefix-two-cycle}
]
#QCIR-G14
forall(1, 2, 3, 4, 5, 6)
forall(7, 8, 9, 10, 11, 12)
exists(13, 14, 15, 16, 17, 18)
exists(19, 20, 21, 22, 23, 24)
exists(25, 26, 27, 28, 29, 30)
forall(31, 32, 33, 34, 35, 36)
output(253)
\end{lstlisting}
Dabei entsprechen die IDs \(1,\ldots,6\) dem Block
\(\mathbf V^1\), die IDs \(7,\ldots,12\) dem Block
\(\mathbf V^2\) und die IDs \(13,\ldots,18\) dem Block
\(\mathbf V^3\). Die Bereiche \(19,\ldots,24\),
\(25,\ldots,30\) und \(31,\ldots,36\) repräsentieren
\(\widehat{\mathbf V}^{1}\),
\(\widehat{\mathbf V}^{2}\) beziehungsweise
\(\widehat{\mathbf V}^{3}\).

Der folgende Auszug zeigt exemplarisch die Schaltkreisübersetzung. Die
ersten fünf Gatter erzwingen für \(a_1\) in \(\mathbf V^1\) genau eines
der drei Labels. Die Gatter \(205\) bis \(211\) kodieren die echte
Inklusion \(I^1<\widehat I^1\). Die abschließenden Gatter setzen die
Antezedens- und Konsequensbausteine zur Gesamtformel zusammen.
\begin{lstlisting}[
    language={},
    caption={Auszug aus der QCIR-Schaltkreismatrix},
    label={lst:qcir-gates-two-cycle}
]
37 = or(1, 3, 5)
38 = or(-1, -3)
39 = or(-1, -5)
40 = or(-5, -3)
41 = and(37, 38, 39, 40)

# I^1 < I^1hat
205 = or(-1, 19)
206 = or(-2, 20)
207 = and(205, 206)
208 = and(19, -1)
209 = and(20, -2)
210 = or(208, 209)
211 = and(207, 210)

# Maximalitaetsbedingungen
226 = or(-148, -211)
227 = or(-176, -218)
228 = or(-204, -225)

# Gleichheit auf dem leeren Z
229 = and()
250 = and()

# Antezedens, Konsequens und Implikation
251 = and(64, 92, 229, 226, 227)
252 = and(120, 239, 249, 250, 228)
253 = or(-251, 252)
\end{lstlisting}
Hier bezeichnen die Gatter \(64\), \(92\) und \(120\) die
\emph{complete}-Kodierungen von
\(\mathbf V^1\), \(\mathbf V^2\) und \(\mathbf V^3\). Die Gatter
\(148\), \(176\) und \(204\) kodieren die entsprechenden
Vergleichslabellings. Die Gatter \(239\) und \(249\) stellen die
Gleichheit auf \(X\) beziehungsweise \(Y\) sicher.

Der Solver wird anschließend ohne weitere Formatkonvertierung
aufgerufen:
\begin{lstlisting}[language={}]
$QUABS zweizyklus.qcir
r UNSAT
\end{lstlisting}
Das Ergebnis stimmt mit der semantischen Betrachtung aus
Beispiel~\ref{ex:preferred-prenex-two-cycle} überein: Das erforderliche
rekombinierende \emph{preferred}-Labelling existiert nicht. Im
verwendeten Auswertungsschema entspricht \texttt{UNSAT} daher dem
Status \texttt{AB}; \texttt{SAT} würde den Status \texttt{UNAB}
bezeichnen.


